\chapter{Inference on a whole scene}\label{ch:inference}

\begin{onething}
To score a real scene, the model runs twice on every tile, with two complementary masks, so that no pixel
is ever predicted while the model can see it; and it runs over hundreds of overlapping tiles, whose
predictions are summed and divided by a count, because a scene is far larger than the 128-pixel patch the
model was trained on.
\end{onething}

Part~V trained a model. You now have a checkpoint: a file of numbers that, loaded into the architecture of
Part~III, turns a masked $128\times128$ patch into a reconstruction. This chapter starts Part~VI, which turns
that ability into an answer an analyst can use. The route has three stages, and this book gives each its
own chapter:
\begin{enumerate}
\item \emph{Inference} (this chapter): run the model over a whole scene and obtain a full-size
  reconstruction in which every pixel was predicted without being seen.
\item \emph{Scoring} (\cref{ch:residualmaha}): turn the difference between the scene and its reconstruction
  into one number per pixel.
\item \emph{Deciding and naming} (\cref{ch:material}): a person picks a threshold, the flagged pixels are
  matched against a library of materials, and the result leaves the system as a georeferenced bundle.
\end{enumerate}
At the end of this chapter you will also walk one real scene through the product's whole chain of actions,
from upload to shortlist, so that you can see where each later chapter fits.

\begin{sloppypar}
Two source files hold the machinery: \file{app/foundation_models/inferencers/segformer_mae_inferencer.py}
(324 lines), which the hyperspectral model inherits unchanged, and
\file{app/foundation_models/inferencers/hyperspectral_segformer_mae_inferencer.py} (105 lines), which
only overrides how the model is built and adds one scoring method. Read them beside this chapter.
\end{sloppypar}

%======================================================================
\section{Two problems inference has that training does not}\label{sec:inference:problems}

\paragraph{Problem 1: every pixel needs a prediction made while it was hidden.}
Recall the rule of \cref{ch:recon}: a pixel must never help reconstruct itself, or the model copies it and
the residual is zero exactly where you hoped to see an anomaly. Training enforces the rule by hiding a random
fraction of the tokens and grading the model only there (\cref{ch:masking}). In training that is enough:
over thousands of patches, every position is hidden many times. At inference you have one particular scene
and need a prediction at \emph{every} pixel of it, each made while that pixel was hidden.

The obvious fix fails in both directions. Run the model once with no mask and every pixel is visible, so
the model copies it. Run it once with a random mask and only the hidden fraction of pixels gets an honest
prediction; the rest have none.

\paragraph{Problem 2: scenes are enormous.}
The model was trained on $128\times128$ patches. A PRISMA scene is about $1210\times1219$
pixels \tagM{docs/lld/frequency-destriper.md}. Again the obvious fixes fail:
\begin{itemize}
\item \emph{Shrink the scene to $128\times128$.} You would average about $9.5\times9.5$ pixels into one, and
  a one-pixel target would vanish before the model ever saw it.
\item \emph{Feed the whole scene in at once.} The encoder's four stages were trained at one scale; a
  $1210\times1219$ input gives a token count at each stage about 90 times larger than anything seen in
  training, and the memory grows with it.
\item \emph{Cut the scene into non-overlapping tiles.} Every tile boundary becomes a seam: a pixel at the
  edge of a tile is predicted from context on one side only, so its residual is systematically worse, and
  the score map shows a grid.
\end{itemize}

The solutions are \emph{two-pass masking} for the first problem and an \emph{overlapping sliding window}
for the second. Before either, the model has to be loaded.

%======================================================================
\section{Loading the model}\label{sec:inference:loading}

The base class \code{FoundationInferencer} does three things in its constructor:
\begin{lstlisting}
self.model = self.build_model().to(self.device)   # 1. build the architecture, empty
self._load_weights(config.checkpoint_path)        # 2. copy the checkpoint's numbers into it
self.model.eval()                                 # 3. switch to inference behaviour
\end{lstlisting}
and \code{_load_weights} reads
\begin{lstlisting}
ckpt = torch.load(path, map_location=self.device, weights_only=False)
self.model.load_state_dict(ckpt["model_state_dict"])
\end{lstlisting}

\paragraph{Why the order matters.}
A \emph{checkpoint}\index{checkpoint} is a dictionary from parameter names (such as
\code{encoder.stages.0.attn.q.weight}) to tensors. It holds numbers, not an architecture. PyTorch cannot
rebuild the network from it; you must build the network first, with exactly the right sizes, and then copy
the numbers into the slots. The copy, \code{load_state_dict}, is \emph{strict} by default: every name in the
file must have a slot of exactly the same shape, and every slot must be filled. That includes the 165-long
normalisation buffers of \cref{ch:bandpipeline} and the spectral compressor sized by $D$
(\cref{ch:compressor}), which is why the resolver reads $D$ from the model's manifest. Build with $D=24$
against a $D=32$ checkpoint and you get an immediate shape error, not a silently wrong score. That
strictness is a feature.

\code{model.eval()} switches off dropout (every unit is kept, \cref{ch:block}) and makes BatchNorm layers
use their stored running statistics instead of the batch's own (\cref{ch:finetune}). Forgetting it gives
reconstructions that change with the batch they happen to be computed in.

%======================================================================
\section{Two-pass masking}\label{sec:inference:twopass}

\subsection{Tokens, keep masks and prediction masks}

Recall three objects from \cref{ch:masking}. The first stage of the encoder cuts a $128\times128$ patch into
non-overlapping $4\times4$ pixel blocks called \emph{tokens}\index{token}; there are
\begin{align*}
\frac{128}{4}\times\frac{128}{4} &= 32\times32\\
&= 1024
\end{align*}
of them, numbered $n=0,\dots,N-1$ with $N=1024$, row by row. A \emph{keep mask}\index{keep mask} is a list of
$N$ zeros and ones: $1$ means ``show this token to the encoder'', $0$ means ``remove it''. The
\emph{prediction mask}\index{prediction mask} is its complement among valid tokens: $1$ means ``this token is
hidden and its prediction will be used''. The \emph{token validity}\index{token validity} $v_n$ is $1$ if the
token's pixels are valid data and $0$ if they are nodata (\cref{ch:masking}).

\subsection{How many passes? The partition argument}\label{sec:inference:partition}

You want each valid token to be hidden in \emph{at least one} pass, and you want to use, for each token,
the prediction from a pass in which it was hidden. Suppose you run $K$ passes and pass $k$ hides the set of
tokens $\mathcal{H}_k$. The cheapest arrangement makes the hidden sets a \emph{partition}\index{partition} of
the tokens: they do not overlap, and together they cover everything,
\begin{equation}
\mathcal{H}_1\cup\mathcal{H}_2\cup\dots\cup\mathcal{H}_K=\{0,1,\dots,N-1\},
\qquad \mathcal{H}_j\cap\mathcal{H}_k=\varnothing\ \text{for } j\neq k.
\label{eq:inference:partition}
\end{equation}
Count the tokens on both sides. Because the sets do not overlap, the count of the union is the sum of the
counts:
\begin{align*}
\abs{\mathcal{H}_1}+\abs{\mathcal{H}_2}+\dots+\abs{\mathcal{H}_K} &= N.
\intertext{If every pass hides the same fraction $\rho$ (the \emph{mask ratio}\index{mask ratio}) of the
tokens, each term is $\rho N$:}
K\rho N &= N.
\intertext{Divide both sides by $N$:}
K\rho &= 1,
\intertext{and so}
K &= \frac{1}{\rho}.
\end{align*}
With $\rho=\tfrac12$ you need $K=2$ passes: hide half, then hide the other half. With $\rho=\tfrac13$ you
would need $K=3$ passes, each showing the model two thirds of its context. The product uses $K=2$,
$\rho=\tfrac12$, always: both of its masking strategies build one mask and its exact complement
(\file{segformer_mae_inferencer.py}, lines 186--216).

\begin{sourcenote}[: the inference mask ratio is not the training mask ratio]
The module docstring says two-pass inference ``matches the training regime (50\% token masking)''
(\file{app/foundation_models/inferencers/segformer_mae_inferencer.py}, line 15). The training-config
default is indeed $0.5$ (\file{app/models/training/training_config.py}, line 81), but the hyperspectral
configurations the courses follow set \code{mask_ratio} to $0.65$
(\file{configs/hyperspectral_segformer_exp_2.json}, line 39, and Indradhanu Heavy's
\file{research/experiments/indradhanu-finetune/FTH3-heavy-joint-more2/config.json}, line 41). So the
shipped models were trained hiding $65\%$ \tagM{configs/hyperspectral_segformer_exp_2.json} of the tokens
and are run at inference hiding $50\%$. A ratio of $0.65$ cannot be used in a two-pass partition at all:
$1/0.65=1.54$ is not a whole number of passes. At inference the model therefore sees \emph{more} context
($50\%$ visible) than in training ($35\%$ visible), which makes its task easier, not harder. The docstring
is out of date; the code is what runs.
\end{sourcenote}

\subsection{The selection is a sum of two masked predictions}

Write $\xhat^{(1)}$ and $\xhat^{(2)}$ for the two passes' reconstructions of a patch, both of shape
\shape{(B, C, 128, 128)}. Let $M^{(1)}$ be a pixel mask that is $1$ on the pixels whose token was hidden in
pass~1, and $M^{(2)}$ the same for pass~2. The code combines them as
\begin{equation}
\xhat = \xhat^{(1)}\odot M^{(1)} + \xhat^{(2)}\odot M^{(2)},
\label{eq:inference:select}
\end{equation}
where $\odot$ is element-wise multiplication (the masks broadcast over the band axis). Because the two hidden
sets partition the tokens, at every pixel exactly one of the masks is $1$:
\begin{align*}
\text{where } M^{(1)}=1,\ M^{(2)}=0:\quad \xhat &= \xhat^{(1)}\cdot1+\xhat^{(2)}\cdot0 = \xhat^{(1)},\\
\text{where } M^{(1)}=0,\ M^{(2)}=1:\quad \xhat &= \xhat^{(1)}\cdot0+\xhat^{(2)}\cdot1 = \xhat^{(2)}.
\end{align*}
So \cref{eq:inference:select} is a \emph{selection}\index{two-pass masking}, not a blend: each pixel takes its value from the one pass
that hid it. Finally the code multiplies by the validity mask, so nodata pixels come out as zero.

\begin{worked}[label={we:inference:select}]{Selecting from two passes on four tokens}
Take a $2\times2$ grid of tokens and one band. Pass~1 hides tokens $(0,0)$ and $(1,1)$; pass~2 hides
$(0,1)$ and $(1,0)$. The two passes return these values (the true values do not matter here):
\[
\xhat^{(1)}=\begin{pmatrix}0.21 & 0.30\\ 0.28 & 0.19\end{pmatrix},\qquad
\xhat^{(2)}=\begin{pmatrix}0.25 & 0.33\\ 0.26 & 0.22\end{pmatrix},\qquad
M^{(1)}=\begin{pmatrix}1&0\\0&1\end{pmatrix},\qquad
M^{(2)}=\begin{pmatrix}0&1\\1&0\end{pmatrix}.
\]
Do it on paper, one entry at a time:
\begin{align*}
\xhat_{0,0} &= 0.21\times1 + 0.25\times0 = 0.21 + 0 = 0.21,\\
\xhat_{0,1} &= 0.30\times0 + 0.33\times1 = 0 + 0.33 = 0.33,\\
\xhat_{1,0} &= 0.28\times0 + 0.26\times1 = 0 + 0.26 = 0.26,\\
\xhat_{1,1} &= 0.19\times1 + 0.22\times0 = 0.19 + 0 = 0.19.
\end{align*}
The result $\begin{pmatrix}0.21&0.33\\0.26&0.19\end{pmatrix}$ uses pass~1 on the diagonal and pass~2 off it.
The value $0.30$ that pass~1 produced at $(0,1)$ is thrown away: in pass~1 that token was visible, so its
prediction may be a copy.
\end{worked}

\subsection{Strategy A: the checkerboard}\label{sec:inference:checker}
\index{checkerboard mask}

The default strategy builds the mask from a checkerboard, in the function \code{checkerboard_token_mask}
of \file{app/foundation_models/components/token_masking.py}:
\begin{lstlisting}
rows = torch.arange(H_tokens) // cell_size          # (32,)
cols = torch.arange(W_tokens) // cell_size          # (32,)
grid = (rows[:, None] + cols[None, :]) % 2          # (32, 32): 1 = visible, 0 = hidden
if invert:
    grid = 1 - grid                                 # pass 2
return grid.flatten().unsqueeze(0).float()          # (1, N) = (1, 1024)
\end{lstlisting}

\paragraph{Broadcasting.}\index{broadcasting} \code{rows[:, None]} has shape \shape{(4, 1)} on a $4\times4$ grid: a column.
\code{cols[None, :]} has shape \shape{(1, 4)}: a row. Adding a column to a row \emph{broadcasts}: the column
is copied across and the row is copied down until both are \shape{(4, 4)}, and then they are added entry by
entry. The result is the addition table.

\begin{worked}[label={we:inference:checker}]{The checkerboard on a $4\times4$ token grid}
With \code{cell_size = 1}, build the table of $\text{row}+\text{col}$:
\[
\begin{array}{c|cccc}
 & \text{col }0 & \text{col }1 & \text{col }2 & \text{col }3\\\hline
\text{row }0 & 0+0=0 & 0+1=1 & 0+2=2 & 0+3=3\\
\text{row }1 & 1+0=1 & 1+1=2 & 1+2=3 & 1+3=4\\
\text{row }2 & 2+0=2 & 2+1=3 & 2+2=4 & 2+3=5\\
\text{row }3 & 3+0=3 & 3+1=4 & 3+2=5 & 3+3=6
\end{array}
\]
Take each entry modulo 2 (the remainder after dividing by 2: $0$ for even, $1$ for odd):
\[
\text{grid}=\begin{pmatrix}0&1&0&1\\1&0&1&0\\0&1&0&1\\1&0&1&0\end{pmatrix}.
\]
In this grid $1$ means \emph{visible}. So pass~1 hides the tokens where $\text{row}+\text{col}$ is
\emph{even}, starting with the top-left token, and pass~2 (\code{invert=True}) hides the odd ones. With
\code{cell_size = 2} the integer division first maps rows $0,1,2,3$ to $0,0,1,1$, and the same steps give
$2\times2$ blocks of tokens.
\end{worked}

\coursefig[0.95\linewidth]{iw-29-two-pass}{Two-pass checkerboard inference (course figure 29.1). Each token's
final value is taken from the pass that hid it, so the combined reconstruction is built entirely from
predictions. See the source check below for which pass hides the top-left token.}{fig:inference:twopass}

\begin{sourcenote}[: which pass hides the top-left token]
The course figure and its caption draw pass~1 hiding the tokens where $\text{row}+\text{col}$ is odd. In the
code, \code{checkerboard_token_mask} returns $1$ for visible and computes $(\text{row}+\text{col})\bmod2$, so
with \code{invert=False} the even tokens are $0$, hidden
(\file{app/foundation_models/components/token_masking.py}, lines 271--284; the comment there reads ``row+col
is even $\to$ 0 (masked)''). The difference is only a relabelling of the two passes and changes no
reconstruction, but if you debug the masks, expect token $(0,0)$ to be hidden in pass~1.
\end{sourcenote}

\subsection{Invalid tokens are always kept}

The keep mask of each pass is built in \code{_checkerboard_keep_mask} as
\begin{lstlisting}
pred_mask = token_validity * (1.0 - checker)    # (B, N): hidden AND valid
keep_mask = 1.0 - pred_mask                     # (B, N): everything else is shown
\end{lstlisting}
Multiplying by \code{token_validity} means a token becomes a prediction target only if it is \emph{both}
valid and selected by the pattern. An invalid token is never removed: it stays in the sequence carrying
the zero signal the model learned to treat as ``no data'', exactly as in training (\cref{ch:masking}).

\begin{worked}[label={we:inference:keep}]{Keep mask for a row of four tokens with one invalid}
Four tokens in one row, the third invalid: $v=(1,1,0,1)$. In pass~1 the checker row is $(0,1,0,1)$ (hide
the first and third). Work entry by entry:
\begin{align*}
1-\text{checker} &= (1-0,\ 1-1,\ 1-0,\ 1-1) = (1,0,1,0),\\
\text{pred} = v\odot(1-\text{checker}) &= (1\cdot1,\ 1\cdot0,\ 0\cdot1,\ 1\cdot0) = (1,0,0,0),\\
\text{keep} = 1-\text{pred} &= (1-1,\ 1-0,\ 1-0,\ 1-0) = (0,1,1,1).
\end{align*}
Only the first token is removed. The third would have been hidden by the pattern, but it is invalid, so it
is kept and nothing is predicted for it.
\end{worked}

\subsection{Strategy B: random halves}

The alternative, \code{masking_strategy="random"}, draws a random half and uses its complement:
\begin{lstlisting}
rand_mask   = (torch.rand(1, N) > 0.5).float()   # (1, N): 1 = visible in pass 1
pred_mask_1 = token_validity * (1.0 - rand_mask)
pred_mask_2 = token_validity * rand_mask          # the exact complement
\end{lstlisting}
Why offer it? A checkerboard is regular. If the model is slightly weaker at token boundaries, and it is,
that weakness lands at the same relative positions in every tile, and a faint grid is printed into the
residual map that an analyst could mistake for structure on the ground. A random mask scatters the
artefacts differently in every tile, and the heavy overlap of tiles (next section) averages them away. The
default in \code{InferenceConfig} is \code{"checkerboard"}
(\file{app/models/training/inference_config.py}, line 96), and the scoring action does not override it, so
the product is deterministic: the same scene and checkpoint give the same score map every time.

\subsection{Cells, blocks and a mismatch to know about}

\begin{sloppypar}
\code{checkerboard_cell_size} (default $1$, \file{inference_config.py}, line 77) sets the size of the
checkerboard cells in tokens. The native-band inferencer used by Indradhanu~V2 is called
\code{SpectralStemInferencer}. It overrides the method \code{_checkerboard_keep_mask} so that a checkpoint
trained with block masking (\code{mask_block} $>1$, \cref{ch:decoderwork}) is hidden in blocks of the same
size. The override is in
\file{app/foundation_models/inferencers/spectral_stem_inferencer.py}, lines 68--77.
\end{sloppypar}

\begin{sourcenote}[: a block-masked checkpoint and the pixel selection]
The override changes which tokens are \emph{hidden}, but \code{infer()} in the parent class builds the
\emph{selection} masks $M^{(1)},M^{(2)}$ of \cref{eq:inference:select} from \code{checkerboard_cell_size},
not from \code{mask_block} (\file{segformer_mae_inferencer.py}, lines 186--199). With \code{mask_block = 2}
and the default cell size $1$ the two patterns disagree. The book's script
\file{research/textbook/tools/compute_inference.py} reproduces both patterns on a $4\times4$ token grid: 8 of
the 16 tokens are taken from a pass in which they were \emph{visible}. The shipped Indradhanu~V2 has
\code{mask_block} $=1$ (the default, \file{app/models/training/training_config.py}, line 199; not set in the
FTH3 config), so the product is unaffected. The block-masked run FT08
(\file{research/experiments/indradhanu-finetune/FT08-gated-blocks/config.json}, line 107: \code{mask_block}
$=2$) was evaluated through the same inferencer with the default cell size
(\file{app/evaluation/models.py}, lines 70--77), so its measured result may include this leak. This was not
verified by rerunning; treat it as an open question for \cref{ch:decoderwork}.
\end{sourcenote}

%======================================================================
\section{The sliding window}\label{sec:inference:sliding}
\index{sliding window}

\subsection{Patch size, stride and overlap}

\code{predict_full_scene} handles a whole scene. Three numbers define its tiling:
\begin{itemize}
\item the \emph{patch size}\index{patch size} $P$, the side of a square tile: $P=128$;
\item the \emph{stride}\index{stride!sliding window} $S$, the distance between the starts of neighbouring
  tiles: the resolver's default for both Indradhanu models is $S=32$
  (\file{backend/allotrope/foundation_models/resolver.py}, lines 160 and 172);
\item the \emph{overlap}\index{overlap} $P-S$, how many pixels two neighbouring tiles share:
  $128-32=96$.
\end{itemize}
If the stride is not given, \code{InferenceConfig} falls back to $P/2=64$
(\file{segformer_mae_inferencer.py}, line 257), which is also the stride used to cut training patches. The
product passes 32: a quarter of the patch, giving heavy overlap. You run inference once per scene, so the
extra compute buys a smoother map.

\subsection{Counting tiles along one axis}\label{sec:inference:count}

The tile starts are decided by \code{PatchPlanGenerator}
(\file{app/utils/patch_generation/generate_patch_plan.py}). Along an axis of length $W$, it places tiles at
$0, S, 2S, \dots$ while they fit, and then places one last tile flush against the far edge. Derive the
count.

A tile starting at $s$ covers pixels $s,\dots,s+P-1$. It fits if its last pixel is inside the axis:
\begin{align*}
s+P-1 &\le W-1.
\intertext{Add $1$ to both sides:}
s+P &\le W.
\intertext{The starts are $s=kS$ for $k=0,1,2,\dots$; substitute:}
kS+P &\le W.
\intertext{Subtract $P$ and divide by $S>0$:}
k &\le \frac{W-P}{S}.
\end{align*}
So the starts that fit are $k=0,1,\dots,\lfloor (W-P)/S\rfloor$, where $\lfloor\cdot\rfloor$ (\emph{floor})
rounds down to a whole number. Counting $0$ as well, that is
\begin{equation}
K_{\text{fit}}=\left\lfloor\frac{W-P}{S}\right\rfloor+1
\end{equation}
tiles. If $(W-P)/S$ is not a whole number, the last fitting tile ends before the edge, leaving a strip
uncovered, so one more tile is placed at $s=W-P$, ``snapped'' back so that it ends exactly at the edge. Then
\begin{equation}
K=\left\lfloor\frac{W-P}{S}\right\rfloor+2=\left\lceil\frac{W-P}{S}\right\rceil+1,
\label{eq:inference:count}
\end{equation}
where $\lceil\cdot\rceil$ (\emph{ceiling}) rounds up. The two forms agree because, for a number that is not
whole, rounding up gives one more than rounding down.

\begin{worked}[label={we:inference:tiles1000}]{Tiles along one axis of a $1000$-pixel scene}
Take $W=1000$, $P=128$, $S=32$. On paper:
\begin{align*}
W-P &= 1000-128 = 872,\\
\frac{W-P}{S} &= \frac{872}{32} = 27.25,\\
\left\lfloor 27.25\right\rfloor &= 27,\\
K_{\text{fit}} &= 27+1 = 28 \quad\text{(starts } 0, 32, \dots, 27\times32=864\text{)}.
\end{align*}
The last fitting tile covers $864,\dots,864+127=991$. Pixels $992$ to $999$ are not yet covered, so a snapped
tile starts at $1000-128=872$. In total
\begin{align*}
K &= 28+1 = 29 = \lceil 27.25\rceil+1 = 28+1.
\end{align*}
The scene has $29\times29=841$ tiles. The book's script, which calls the product's own
\code{PatchPlanGenerator}, gives $29$ per axis and $841$ tiles \tagC.
\end{worked}

\begin{worked}[label={we:inference:tilesprisma}]{A real PRISMA frame, $1210\times1219$}
A PRISMA L2D frame of $1210\times1219$ pixels \tagM{docs/lld/frequency-destriper.md}. Rows first:
\begin{align*}
\frac{1210-128}{32} &= \frac{1082}{32} = 33.8125,\\
K_{\text{rows}} &= \lceil 33.8125\rceil+1 = 34+1 = 35.
\end{align*}
Columns:
\begin{align*}
\frac{1219-128}{32} &= \frac{1091}{32} = 34.09375,\\
K_{\text{cols}} &= \lceil 34.09375\rceil+1 = 35+1 = 36.
\end{align*}
Tiles:
\begin{align*}
35\times36 &= 1260.
\end{align*}
The model takes \code{inference_batch_size} tiles at a time. The scoring action passes the resolver's
default of $8$ (\file{resolver.py}, line 62), so the number of batches is
\begin{align*}
\left\lceil\frac{1260}{8}\right\rceil = \lceil 157.5\rceil = 158,
\end{align*}
and since every batch runs the network twice (two passes),
\begin{align*}
158\times2 = 316 \text{ forward passes.}
\end{align*}
The script confirms $35$, $36$, $1260$ tiles and $316$ forward passes \tagC.
\end{worked}

\begin{sourcenote}[: when $(W-P)/S$ is a whole number, the code adds a duplicate tile]
\Cref{eq:inference:count} is what the generator's docstring intends, and it holds whenever $(W-P)/S$ is not
a whole number. When it \emph{is} whole, the last fitting tile already ends at the edge, yet the loop in
\file{generate_patch_plan.py} (lines 66--80) appends it, steps on by $S$, finds the next tile overhanging and
appends the snapped start $W-P$, which is the same start again. For $W=1024$: $(1024-128)/32=28$, the formula
gives $29$ tiles, and the code produces 30 starts of which 29 are distinct; for $W=128$ it produces the start
$0$ twice \tagC. The duplicate tile wastes one tile's compute per row or column and counts that tile twice in
the average below. With the deterministic checkerboard both copies are identical, so the reconstruction is
unchanged; with random masking the last tile gets double weight. PRISMA frames are not multiples of 32
plus 128, so the product rarely meets the case.
\end{sourcenote}

\coursefig[0.95\linewidth]{iw-29-sliding}{Sliding-window inference (course figure 29.2). $128\times128$ tiles
start every 32 pixels, so a $1000\times1000$ scene needs $29\times29=841$ tiles; an interior pixel lies in
$128/32=4$ tiles along each axis and so in $4\times4=16$ tiles in all.}{fig:inference:sliding}

\subsection{How many tiles cover each pixel}\label{sec:inference:coverage}
\index{coverage count}

Along one axis, a pixel at position $x$ lies in every tile whose start $s$ satisfies $s\le x\le s+P-1$. Away
from the edges the starts are $S$ apart, so the number of such starts is the length $P$ of the window
divided by the gap $S$ between starts:
\begin{align*}
\frac{P}{S} &= \frac{128}{32} = 4.
\end{align*}
In two dimensions a pixel lies in a tile only if both its row and its column do, so the counts multiply:
\begin{align*}
4\times4 &= 16 \text{ tiles.}
\end{align*}
Sixteen separate reconstructions of every interior pixel, each made with the tile in a different position
and so with different context visible around it. Near the edges the count drops: a pixel in the first 32
columns is in only one tile along that axis. Near the far edge something less obvious happens: the snapped
tile overlaps its neighbours by more than $P-S$, so a strip of pixels is covered \emph{five} times.
\Cref{fig:inference:coverage} plots the count along one axis of the $1000$-pixel scene, from the script's
output \tagC.

\begin{figure}[htbp]\centering
\begin{tikzpicture}
\begin{axis}[width=0.95\linewidth,height=5.2cm,xmin=0,xmax=1000,ymin=0,ymax=6.4,
  xlabel={column $x$ (pixels)},ylabel={tiles containing $x$},ytick={0,1,2,3,4,5},
  xtick={0,96,200,400,600,800,872,1000},grid=major,grid style={gray!20}]
\addplot[const plot,very thick,spat] coordinates {(0,1)(32,2)(64,3)(96,4)(872,5)(896,4)(928,3)(960,2)(992,1)(1000,1)};
\node[font=\small,anchor=south] at (axis cs:480,4.05) {interior: $128/32=4$};
\node[font=\small,anchor=south,loss] at (axis cs:884,5.05) {snapped tile: 5};
\end{axis}
\end{tikzpicture}
\caption[Tiles covering each column of a 1000-pixel axis]{Number of tiles containing each column of a $1000$-pixel axis with $P=128$, $S=32$, computed with
the product's \texttt{PatchPlanGenerator} (\texttt{research/textbook/tools/compute\_inference\_output.txt}).
The two-dimensional count at a pixel is the product of its row count and its column count: 16 in the
interior, 1 at a corner, 25 where the two snapped strips cross.}\label{fig:inference:coverage}
\end{figure}

%======================================================================
\section{Accumulate, then average}\label{sec:inference:average}

\subsection{Two scene-sized accumulators}

Storing 841 or 1260 reconstructed tiles and averaging them at the end would need hundreds of copies of the
scene. Instead, \code{predict_full_scene} keeps two running totals the size of the scene:
\begin{lstlisting}
recon_sum = torch.zeros(c, h, w)    # (C, H, W): sum of the predictions that reached each pixel
count     = torch.zeros(1, h, w)    # (1, H, W): how many predictions reached each pixel
...
for j, idx in enumerate(valid_indices):
    r, cs = batch_coords[idx]
    patch_eroded = eroded_mask[:, r:r+ps, cs:cs+ps]          # (1, 128, 128), 0 or 1
    recon_sum[:, r:r+ps, cs:cs+ps] += recon[j] * patch_eroded  # (C, 128, 128)
    count    [:, r:r+ps, cs:cs+ps] += patch_eroded
reconstruction = torch.where(count > 0, recon_sum / count, scene)   # (C, H, W)
\end{lstlisting}
In symbols, write $\xhat_t(p)$ for tile $t$'s prediction at pixel $p$ and $w_t(p)\in\{0,1\}$ for the eroded
validity of $p$ (defined below; $0$ if $p$ is not in tile $t$). Then
\begin{equation}
\hat{x}(p)=
\begin{cases}
\dfrac{\sum_t w_t(p)\,\xhat_t(p)}{\sum_t w_t(p)} & \text{if } \sum_t w_t(p)>0,\\[2ex]
x(p) & \text{if } \sum_t w_t(p)=0.
\end{cases}
\label{eq:inference:average}
\end{equation}
The top line is a \emph{weighted average}\index{weighted average}: the weights are $0$ or $1$, so it is the
plain mean of the predictions that count.

\subsection{Three details}

\paragraph{The weights come from an eroded mask.}
\emph{Erosion}\index{erosion} shrinks the valid region: a pixel survives only if every pixel in a
$k\times k$ square around it is valid (\cref{ch:masking}). Inference uses $k=15$
(\file{inference_config.py}, line 109), so any pixel within $\lfloor15/2\rfloor=7$ pixels of nodata gets
weight $0$. Such pixels are predicted from context that is partly nodata, and their reconstructions are
unreliable. With weight $0$ they add nothing to the sum \emph{and} nothing to the count: not a small
amount, nothing at all.

\paragraph{Dividing by the count makes the average right everywhere.}
An interior pixel has count 16, a corner pixel count 1, a pixel in the snapped strip up to 25. Dividing by
the actual count gives a true mean in every case with no special handling of edges.

\paragraph{Where nothing arrived, the scene stands in for itself.}
If no tile delivered a usable prediction (count $0$), the reconstruction is set to the input. The residual
there is $x(p)-x(p)=0$, so the pixel scores zero and can never be flagged. Leaving zeros or NaN in such
pixels would have produced huge or undefined residuals instead.

\begin{worked}[label={we:inference:avg}]{Averaging three tiles, one of them eroded}
A pixel $p$ in one band lies in three tiles. Their predictions and weights are
\[
\xhat_1(p)=0.24,\ w_1=1;\qquad \xhat_2(p)=0.28,\ w_2=1;\qquad \xhat_3(p)=0.90,\ w_3=0.
\]
(Tile~3 is the one in which $p$ sits within 7 pixels of a nodata hole, and its prediction is spoiled.) On
paper:
\begin{align*}
\text{sum} &= 1\times0.24 + 1\times0.28 + 0\times0.90\\
&= 0.24 + 0.28 + 0\\
&= 0.52,\\
\text{count} &= 1+1+0 = 2,\\
\hat{x}(p) &= \frac{0.52}{2} = 0.26.
\end{align*}
Had the third tile been given a small weight such as $0.1$ instead of $0$, the sum would be
$0.52+0.09=0.61$, the count $2.1$, and the average $0.61/2.1=0.290$: a spoiled prediction still pulls
the answer by $0.03$.
\end{worked}

\paragraph{Why averaging smooths the map.}
If the errors of the 16 predictions of an interior pixel were independent with equal variance $\sigma^2$,
the variance of their mean would be
\begin{align*}
\operatorname{Var}\!\left(\frac{1}{16}\sum_{t=1}^{16}\xhat_t\right)
&= \frac{1}{16^2}\sum_{t=1}^{16}\operatorname{Var}(\xhat_t)
 = \frac{16\,\sigma^2}{256}
 = \frac{\sigma^2}{16},
\end{align*}
a standard deviation four times smaller. In practice the predictions share most of their context and are
far from independent \tagA, so the real gain is smaller, but the direction is right: overlapping tiles
average away the noise that any one tile position adds. A smooth reconstruction is expected; a crisp one
would suggest the model had found a way to copy.

\subsection{The tile validity filter}

A tile is only processed if at least $10\%$ of its pixels are valid:
\begin{lstlisting}
MIN_VALID_FRACTION = 0.1
valid_frac = m.float().mean().item()      # fraction of the tile's pixels that are valid
if valid_frac >= MIN_VALID_FRACTION:
    valid_indices.append(i)
\end{lstlisting}
Tiles that are more than $90\%$ nodata are skipped: the model never saw such patches in training and its
output on them is meaningless.

\begin{sourcenote}[: the comment says 0.4, the code says 0.1]
The comment above this line says the threshold ``matches training: patches with $<40\%$ valid pixels were
discarded'' (\file{segformer_mae_inferencer.py}, lines 280--284). The training threshold is indeed $0.4$
(\code{MIN_VALID_PIXEL_FRACTION}, \file{app/foundation_models/trainers/hyperspectral_segformer_mae_trainer.py},
line 34), but inference uses $0.1$. The intent, skipping mostly empty tiles, is the same; the numbers
differ. The course (IW29) already flags this.
\end{sourcenote}

\subsection{Memory, roughly}

One full cube of $165$ bands at $1000\times1000$ pixels in 4-byte \code{float32} holds
\begin{align*}
165\times1000\times1000 &= 165{,}000{,}000 \text{ numbers},\\
165{,}000{,}000\times4 &= 660{,}000{,}000 \text{ bytes} = 660\text{ MB}.
\end{align*}
Inference keeps at least three such arrays: the scene, \code{recon_sum}, and the final reconstruction,
about $3\times660=1980$~MB, about 2~GB, before any network activations. For the $1210\times1219$ frame each
array is $973{,}493{,}400$ bytes \tagC, close to 1~GB, so about 3~GB in total. That is why
\code{inference_batch_size} is exposed as a per-action override with a conservative default of 8: the
resolver's own comment calls it ``the main memory dial when running on a CPU-only Mac VM''
(\file{resolver.py}, lines 59--62).

\begin{sourcenote}[: two defaults for the batch size]
\code{InferenceConfig.inference_batch_size} defaults to 16 (\file{inference_config.py}, line 104), but the
scoring action always passes the resolver's \code{default_batch_size} of 8 unless the user overrides it
(\file{backend/allotrope/action_types/_anomaly_scoring_run.py}, line 303). In the product the batch size is 8.
\end{sourcenote}

%======================================================================
\section{Shapes from scene to reconstruction}\label{sec:inference:shapes}

Every tensor on the inference path, for the $1210\times1219$ PRISMA frame and the 165-band Indradhanu:
\begin{lstlisting}
scene            (165, 1210, 1219)  float32, from filtered_vendable.pkl
mask             (1, 1210, 1219)    1 = spatially valid (any band valid)
eroded_mask      (1, 1210, 1219)    after a 15 x 15 erosion
coords           1260 pairs (r, c)  35 row starts x 36 column starts
patches          (8, 165, 128, 128) one batch of tiles
patch_masks      (8, 1, 128, 128)
token_validity   (8, 1024)          32 x 32 tokens per tile
keep_mask_1/_2   (8, 1024)          1 = show, 0 = hide (complements among valid tokens)
pass1_pixels     (1, 1, 128, 128)   token mask expanded to pixels (nearest-neighbour)
x_hat_1, x_hat_2 (8, 165, 128, 128) the two passes
recon            (8, 165, 128, 128) x_hat_1*pass1_pixels + x_hat_2*pass2_pixels, times mask
recon_sum        (165, 1210, 1219)  running sum
count            (1, 1210, 1219)    running count, 0..25
reconstruction   (165, 1210, 1219)  recon_sum / count, or the scene where count = 0
\end{lstlisting}

%======================================================================
\section{From reconstruction to score map}\label{sec:inference:toscore}

\subsection{The scoring call}

The action does not use the inferencer's own scoring method. It calls one shared numpy function,
\code{compute_score} in \file{app/utils/anomaly_detection/scoring.py}:
\begin{lstlisting}
score = compute_score(cube, reconstruction, eroded_keep_mask,   # (C,H,W), (C,H,W), (H,W)
                      method=method, combined_weight=sam_l1_alpha)   # -> (H, W) float32
\end{lstlisting}
It returns one number per pixel, higher meaning more anomalous, and sets every pixel outside the keep mask to
exactly $0$ (\file{scoring.py}, line 114). Five methods exist. Four of them measure the \emph{size} of the
residual $\vect{r}=\vect{x}-\xhat$:
\begin{itemize}
\item \code{L1}: the mean of $\abs{r_c}$ over the bands;
\item \code{MSE}: the mean of $r_c^2$ (not offered for Indradhanu);
\item \code{SAM}: the spectral angle between $\vect{x}$ and $\xhat$ in radians (\cref{ch:losses});
\item \code{combined}: $w\cdot\text{L1}_{\text{norm}}+(1-w)\cdot\text{SAM}_{\text{norm}}$, each divided by its
  maximum over valid pixels, $w=0.5$ by default.
\end{itemize}
The fifth, \code{residual_mahalanobis}, measures how \emph{unusual} the residual's direction is, and has
been Indradhanu's default since 2026-10-15 (\file{resolver.py}, lines 152--158). \Cref{ch:residualmaha} is
devoted to it. The inferencer's own \code{compute_anomaly_scores}, which returns a dictionary of L1 and SAM
maps, exists for notebooks; the product never calls it.

\coursefig[0.9\linewidth]{iw-30-scoring}{The \texttt{combined} scoring method (course figure 30.1): per-pixel L1
and SAM maps are each divided by their maximum over the valid pixels and blended with weight 0.5. Since
2026-10-15 this is no longer Indradhanu's default.}{fig:inference:combined}

\begin{worked}[label={we:inference:combined}]{The combined score and its fragile denominator}
L1 is in reflectance units and SAM in radians; adding them directly would be adding metres to kilograms. So
each is divided by its own maximum over the scene's valid pixels. Suppose the scene maxima are
$\max\text{L1}=0.05$ and $\max\text{SAM}=0.40$~rad, and one pixel has $\text{L1}=0.025$ and
$\text{SAM}=0.10$~rad \tagA. Normalise:
\begin{align*}
\text{L1}_{\text{norm}} &= \frac{0.025}{0.05} = 0.50,\\
\text{SAM}_{\text{norm}} &= \frac{0.10}{0.40} = 0.25.
\end{align*}
Blend with $w=0.5$:
\begin{align*}
\text{score} &= 0.5\times0.50 + 0.5\times0.25\\
&= 0.250 + 0.125\\
&= 0.375.
\end{align*}
Now let one dead detector element produce a single pixel with $\text{L1}=0.50$, ten times the old maximum.
The same pixel becomes
\begin{align*}
\text{L1}_{\text{norm}} &= \frac{0.025}{0.50} = 0.05,\\
\text{score} &= 0.5\times0.05 + 0.5\times0.25 = 0.025 + 0.125 = 0.150.
\end{align*}
One outlier has compressed every other pixel's L1 term by a factor of ten. The symptom is a dark score map
with one bright dot; the check is the percentiles in \code{stats.json}, written for every model.
\end{worked}

\subsection{What the scoring action writes}

For each model, in \file{models/<codename>/} under the action's output folder, where the folder name is the
codename in lower case with spaces turned into underscores (\code{_safe_dirname}):

\begin{center}\small
\begin{tabular}{@{}lll@{}}
\toprule
file & shape and type & contents\\
\midrule
\texttt{anomaly\_score.tif} & $(H,W)$ float32 & the score map\\
\texttt{anomaly\_score.png} & image & preview; square-root stretch for squared distances\\
\texttt{reconstruction.tif} & $(C,H,W)$ float32 & the full-scene reconstruction\\
\texttt{reconstruction.png} & image & RGB preview of the reconstruction\\
\texttt{stats.json} & JSON & method, patch, stride, timings, score percentiles\\
\bottomrule
\end{tabular}
\end{center}

The preview stretch is a square root for the classical detectors \emph{and} for \code{residual_mahalanobis},
because both produce squared distances with a heavy right tail; linear otherwise
(\file{_anomaly_scoring_run.py}, line 511). Every GeoTIFF is written with an identity transform
(\file{_anomaly_scoring_render.py}, line 61): it carries no geography until export (\cref{ch:material}).

%======================================================================
\section{One scene through the whole product}\label{sec:inference:chain}

The components of this book stop being separate here. Follow one PRISMA scene from upload to a shortlist;
each step names the chapter that explained it. The numbers are those of the anomaly course's end-to-end
walk (AD13): a scene of $1210\times1219$ pixels, with the fractions lost at each step assumed for
illustration \tagA.

\begin{figure}[htbp]\centering
\begin{tikzpicture}[node distance=4mm,
  box/.style={draw,rounded corners,fill=#1!8,draw=#1!70!black,minimum width=37mm,minimum height=8mm,align=center,font=\small},
  lab/.style={font=\scriptsize,align=left,anchor=west}]
\node[box=maskgrey] (up) {upload + \texttt{scene\_onboard}};
\node[box=maskgrey,below=of up] (pr) {Project (bound to the Scene)};
\node[box=spec,below=of pr] (bf) {\texttt{band\_filter\_apply}};
\node[box=spec,below=of bf] (ss) {\texttt{scene\_segmentation}};
\node[box=attn,below=of ss] (as) {\texttt{anomaly\_scoring}};
\node[box=loss,below=of as] (dp) {\texttt{anomaly\_detection\_prep}};
\node[box=spec,below=of dp] (sm) {\texttt{spectral\_library\_match}};
\node[box=maskgrey,below=of sm] (ex) {\texttt{POST /actions/\{id\}/export}};
\foreach \a/\b in {up/pr,pr/bf,bf/ss,ss/as,as/dp,dp/sm,sm/ex} \draw[-Stealth] (\a) -- (\b);
\node[lab,right=5mm of up] {vendable: cube $(165,1210,1219)$, validity; no coordinates};
\node[lab,right=5mm of pr] {nothing computes};
\node[lab,right=5mm of bf] {\texttt{filtered\_vendable.pkl}, nearest-valid fill};
\node[lab,right=5mm of ss] {\texttt{keep\_mask.tif}: water, cloud, shadow, vegetation out};
\node[lab,right=5mm of as] {two passes, 1260 tiles, average; one score map per model};
\node[lab,right=5mm of dp] {composite score; stops at \texttt{needs\_threshold}};
\node[lab,right=5mm of sm] {native cube; \texttt{family\_map.tif}, \texttt{match\_map.tif}};
\node[lab,right=5mm of ex] {re-read raw file for CRS; zip of GeoTIFF, Shapefile, CSV};
\end{tikzpicture}
\caption{The product's action chain for a hyperspectral scene (AD13). Teal steps transform spectra, plum runs
the model, amber is where a person decides.}\label{fig:inference:chain}
\end{figure}

\begin{enumerate}
\item \textbf{Upload and onboarding} (\cref{ch:product}). The api streams the file to a staging folder and
  queues a \code{scene_onboard} job; it does no analysis itself. The worker builds the vendable with the
  eight-stage band pipeline (\cref{ch:bandpipeline}): a cube \shape{(165, 1210, 1219)} of reflectance, a
  validity cube of the same shape, and 165 wavelengths. Note what is absent: no latitude, no longitude, no
  projection.
\item \textbf{Project} (\cref{ch:product}). Bound to the scene. Nothing computes.
\item \textbf{\texttt{band\_filter\_apply}}. Re-vends the scene with explicit band settings and writes
  \file{filtered_vendable.pkl}, filling each invalid pixel with its nearest valid neighbour so the patch
  embedding never meets a cliff (\cref{ch:bandpipeline}).
\item \textbf{\texttt{scene\_segmentation}}. Thresholds spectral indices into masks and writes
  \file{keep_mask.tif}, the region worth analysing.
\item \textbf{\texttt{anomaly\_scoring}}. This chapter and the next: tile, two passes, average, score. A
  classical detector such as MNF-RX (\cref{ch:localrx}) can run in the same action; its covariance is
  estimated only on the keep mask.
\item \textbf{\texttt{anomaly\_detection\_prep}}. Combines the models' maps into one composite and
  \emph{stops} at the status \code{needs_threshold} until a person chooses a percentile
  (\cref{ch:material}).
\item \textbf{\texttt{spectral\_library\_match}}. Hyperspectral only. Names the likely material family of each
  flagged pixel, reading the native onboarding cube, not the filtered one (\cref{ch:material}).
\item \textbf{Export}. Coordinates finally appear: the raw file is re-read for its georeferencing, and a zip
  leaves the system (\cref{ch:material}).
\end{enumerate}

\begin{worked}[label={we:inference:funnel}]{The funnel from 1.47 million pixels to a shortlist}
Assume onboarding loses $6\%$ of the pixels to cloud and swath edges, segmentation keeps $69\%$ of the
rest, and the analyst flags the top $0.5\%$ of the kept pixels \tagA. On paper:
\begin{align*}
\text{pixels} &= 1210\times1219 = 1{,}474{,}990,\\
\text{valid} &= 1{,}474{,}990\times0.94 = 1{,}386{,}490.6 \approx 1{,}386{,}491,\\
\text{kept} &= 1{,}386{,}491\times0.69 = 956{,}678.8 \approx 956{,}679,\\
\text{flagged} &\approx 956{,}679\times0.005 = 4{,}783.4.
\end{align*}
The reduction factor is
\begin{align*}
\frac{1{,}474{,}990}{4{,}783} &\approx 308.4.
\end{align*}
From about 1.47 million pixels to about 4,800: the system narrows, and a person reviews the shortlist.
\Cref{ch:material} shows that the exact count the code flags is $4{,}784$, and that with a model-only chain
the count can be quite different, because of which pixels the percentile is taken over.
\end{worked}

\paragraph{What can go wrong at each step.} The anomaly course lists the failures seen in practice:
the wrong sensor selected (a \code{ValueError}, job failed); a class renamed in \file{vendables.py}
(stored pickles unreadable); segmentation thresholds wrong for the terrain (the keep mask removes the
target); the worker crash-looping (jobs sit in \code{queued} while the api reports healthy, which has
happened twice, see CLAUDE.md); normalisation statistics wrong for the sensor (every pixel scores high); a
threshold too tight (the real anomaly ranked a little below the cut); and georeferencing that cannot be
resolved (\code{422 crs_missing}).

%======================================================================
\begin{confusions}
\confusion{Does the scene get resized to 128 by 128?}{No. It is cut into overlapping $128\times128$ tiles,
each reconstructed separately at full resolution, and the predictions are averaged back into a full-size
cube.}
\confusion{Are the two passes twice as slow?}{Yes, exactly twice the network calls. It is the price of the
rule that no pixel predicts itself, and it is not negotiable.}
\confusion{Is inference masking the same as training masking?}{No. Training hides a random $65\%$ and
grades only there; inference hides exactly half, twice, with complementary masks, and uses every hidden
prediction. The partition argument forces $\rho=1/K$.}
\confusion{Why does a stride of 32 cost sixteen times more than a stride of 128?}{The number of tiles along
each axis grows about four times ($P/S=4$ instead of $1$), and the two axes multiply: $4\times4=16$ times as
many tiles, each giving every pixel one more prediction.}
\confusion{Pixels near a nodata hole score zero. Is the model ignoring them?}{If no tile delivered a usable
prediction for a pixel (all its weights eroded to 0), the reconstruction falls back to the input and the
residual is exactly zero. If some did, the pixel is scored normally from those. Erosion of 15 removes only
pixels within 7 of nodata from the average.}
\confusion{The score map has a faint grid.}{Checkerboard artefacts at token boundaries. The random strategy
or a different stride changes them; the overlap of 16 tiles already damps them.}
\end{confusions}

\begin{checkyourself}
\item\label{ex:inference:1} Explain in two sentences why one pass with a $50\%$ mask is not enough, and why
  one pass with no mask is worse.
\item\label{ex:inference:2} A scene is $1200\times900$ pixels, $P=128$, $S=32$. Compute the number of tile
  starts along each axis with \cref{eq:inference:count}, the number of tiles, the number of batches of 8
  and the number of forward passes. Show every step.
\item\label{ex:inference:3} With \code{cell_size = 2}, write out the $4\times4$ token grid that pass~1 hides,
  and the one pass~2 hides.
\item\label{ex:inference:4} A pixel lies in four tiles with predictions $0.31, 0.29, 0.30, 0.70$; the last
  tile's weight is $0$ because of erosion. Compute the averaged reconstruction, and the residual if the
  observed value is $0.36$. What would the residual be if all four weights were $0$?
\item\label{ex:inference:5} Along a $W=600$ axis with $P=128$, $S=32$, list the last three tile starts and
  give the coverage count (tiles containing it) of column $500$.
\item\label{ex:inference:6} Why must a checkpoint be loaded into a model built with the same $D$, and what
  happens if it is not?
\end{checkyourself}

\nextchapter{ch:residualmaha}
